\section{怎样分析句子的结构}

句子是由词或词组构成的能够表达一个完整意思的语言单位。根据句子的结构不同，句子分为单句与复句。如果我们能够准确地分析单句的结构，那么分析具有两个以上单句（叫做分句）的复句，也不在话下了，因为分析复句的分句结构同分析单句的结构的方法是一样的。

单句分单句和双句。单句包括独词句（指用一个名词或名词性词组构成的句子，如“火！” “山！”“蔚兰的大海。”）和无主句（只有谓语部分。没有主语部分的句子，如“静！”“严禁吸烟！”“爱护公物。”）；双句包括主谓句（由主语和谓语两部分构成的句子，如
“\uuline{人民}\underline{万岁}！”
“\uuline{我}\uline{热爱\\祖国}”）
和省略句（指由于承前或蒙后而省略了主语的句子，如
“\uuline{我}\underline{爱高山}，
（\uuline{我}）\underline{也爱大海}。”）。

有的中学生说，对于“社会主义好”，“人民热爱共产党”这一类结构简单的单句，我们会分析；对于结构复杂的单句就不知从哪儿下手了。比如，秦牧的散文《土地》中有这么一个长句子：“当你坐在飞机上，看着我们无边无际的景象覆盖上一张绿色地毯的大地的时候，当你坐在汽车上，倚着车窗看万里平畴的时候，或者，在农村里，看到一个老农捧起一把泥土，仔细端详，想鉴定它究竟适宜于种植什么谷物和蔬菜的时候，或者，当你自己随着大伙在田里插秧，黑油油的泥土吱吱地冒出脚指缝的时候，你曾否为土地涌现过许许多多的遐想——想起它的过去，它的未来，想起世世代代的劳动人民为要成为土地的主人，怎样斗争和流血，想起在绵长的历史中，我们每一块土地上面曾经出现过的人物和事迹，他们的痛苦、怨恨、希望、期待的心情？”（见高中语文课本第二册）究竟该如何分析呢？

其实，不管句子结构如何复杂，我们都可以从抓住主干入手，化繁为简，逐层剖析，弄清句子的结构。具体说来，可按下面的办法做。

1. 找寻主语，弄清句子类型。

我们知道，除了独词句（如“海！”、“多好的阳光！”等）、无主句（如“冲！”“禁止攀折花木！”“出太阳了。”等）这两种不是主语和谓语全都具有的“单部句”外，语言中最多最常见的一类句子是主语、谓语全部具备的“双部句”（在一定语言环境里临时省略了主语的“省略句”，因为省略的主语是确定的，可以补出来，因此仍然属于“双部句”）。我们在分析句子的结构时，首先要全力弄清整个句子有没有主语，有几个主语，以确定这个句子究竟是单句，还是复句，是主谓完全具备的主谓句，还是没有主语只有谓语的非主谓句。例如上面所引的《土地》一文中的长句，我们稍加分析就可判定：全句只有一个被陈述的对象--主语，即“你曾否为土地涌现过许许多多的遐想”中的“你”，其他所有的词语都没有资格充当这个长句的被陈述的对象。至于“你”怎么样呢？紧接着就作了陈述。据此可以判断这是一个单句中的“双部句”，当然也是一个主谓具备的“主谓句”。

再如：“在旧社会，多少从事科学文化事业的人们，向往着国家昌盛、民族复兴、科学文化繁荣啊！”稍加分析，就可找出全句被陈述的对象--主语“人们”，而“人们”怎么样呢？接着就进行了陈述，说明这是一个单句中的“双部句”，也是一个主谓具备的“主谓句”。

2. 切分主谓，明确句子两大部分。

我们知道，主语和谓语是构成主谓句的主要成分，主谓句是句子的典型。主语是谓语陈述的对象，谓语是对主语加以陈述的。一般说来，在汉语里，主语出现在谓语的前边。因此，当找出句子的主语后，我们就可以在主语后面画两条竖线，将全句一分为二。这样，一个句子的主语部分和谓语部分就明确起来了。例如上面所引《土地》中的长句，我们可以在句中的“你”之后画上“$\Vert$”，表示两条竖线后面的部分，就是对主语部分加以陈述的谓语部分。又如：

大多数的学生和教师 $\Vert$，都是把为人民服务、为革命事业献身作为自己的最高理想。

\uuline{清理古代文化的发展过程，剔除其封建性的糟粕，\\
吸收其民主性的精华} $\Vert$，是发展民族新文化提高民族自信心的必要条件。

3. 找出中心词，突现句子基本结构。

当明确了主语部分和谓语部分之后，就要努力找出主语部分的中心词--主语和谓语部分的中心词--谓语，以及动词谓语涉及的对象--宾语。这样，一个句子的基本结构就突现出来了。当然，上面所引的《土地》中的长句，它的主语部分简单，就是一个“你”字；而一般的句子，在主语部分的中心词之前往往有修饰或限制这个中心词的定语，甚
至是非常复杂的定语，这就要细心辨析。例如：
\newpage
①多少从事科学文化事业的\uline{人们} $\Vert$，\uline{向往}着······。

（主语部分中心词“人们”之前，有“多少”“从事科学文化事业”作为修饰或限制这个主语的定语。）

②大多数的\uuline{学生和教师} $\Vert$，都是把······\uline{作为}······理想。
（主语部分中心词是“学生和教师”这个联合词组，在它之前，有“大多数”作为限制这个主语的定语。）

③\uuline{清理······过程，剔除······糟粕，吸收······精华} $\Vert$，\uline{是}······条件。

（这个句子的主语是由三个比较复杂的动宾词组所构成的联合词组来充当的，作为全句被陈述的对象。）

至于谓语部分的中心词是指陈述前面主语的主要动词或形容词。一般句子中只有一个谓语中心词或一个并列的谓语中心词组，象刚才所说的三个例句中的“向往”、“作为”、“是”。而在“连动式”或者“兼语式”谓语部分中，可以出现两个或两个以上的谓语部分中心词。例如（
\uuline{\ \ \ }表示主语；
\uline{\ \ \ }表示谓语；
\uwave{\ \ \ }表示宾语；
表示兼语）：

\uuline{黄老师} $\Vert$ \uline{站起来}\ \uline{迎接}\uwave{他们}。（连动式）
这件\uuline{事} $\Vert$ \uline{使}我非常\uline{着急}。（兼语式）

另外，汉语中很多动词要带宾语，如果这个动词作为句子的谓语部分中心词，那么它所支配的对象宾语也是句子的基本成分之一。例如上面所说的三个例句，其基本结构是：

①······\uuline{人们} $\Vert$ \uline{向往}着\uwave{国家昌盛、民族复兴、科学文化\\繁荣}啊！（三个并列的主谓词组充当谓语中心词“向往”的宾语。）

②······\uuline{学生和教师} $\Vert$，都是把······\uline{作为}······\uwave{理想}。

③\uuline{清理······过程，剔除······糟粕，吸收······}\ \uline{精华} $\Vert$，\\
\uline{是}······\uwave{条件}。

\newpage

当我们注意到以上所说的这些问题之后，就可以顺利地确定一个句子的基本结构。比如《土地》中的那个长句的基本结构就是：

\uuline{你} $\Vert$ 曾否为土地\uline{涌现}过许许多多的\uwave{遐想}？

4.根据主要成分，分析附加成分。

我们只要确定了句子的主要成分主语、谓语和宾语之后，附加成分的分析也就迎刃而解了。因为一般说来，定语放在主语或宾语之前，修饰后边的名词，定语之后常常用结构助词“的”，状语放在谓语之前，修饰后边的动词，形容词，状语之后常常用结构助词“地”；补语放在谓语之后，补充说明前面的动词、形容词，补语的前边常常用结构助词“得”。有的句子虽然不用“的、地、得”作为定语、状语和补语的标志，但是从这些附加成分与主要成分之间的关系和位置也是可以看出来的。

另外，在不少句子中，为了强调或把太长的状语分出来好念，经常将某些副词或介词结构充当的状语，提到全句的前面，这一点应该注意，千万不要误将介词结构看成主谓，因为介词结构是不能充当主语成分的。

例如（用（）表示定语，[ ]表示状语，（>）表示补语）：

①[从很早的古代起]，（我们中华民族）的\uuline{祖先} $\Vert$ [就]
\uline{劳动、生息、繁殖}(在这块广大的土地上)。

（“从······古代起”
“在······土地上”
是介词结构，分别作句子的状语或补语；加点的“的”是定语结构助词，不能误将“我们”作为句子的主语。）

②（祖国）的（更加美好）的\uuline{未来} $\Vert$，[越来越鲜明]地\uline{出现} $\langle$在我们的眼前$\rangle$。

（“在······眼前”是介词结构，作“出现”的补语；“的”是定语结构助词，不能误把“祖国”作为句子的主语；“地”是状语结构助词。）

③（英雄）的\uuline{火箭炮手} $\Vert$ [把越寇]\uline{打}得$\langle$ 焦头烂额 $\rangle$。

（“把越寇”是介词结构，充当“打”的状语；“的”是定语结构助词，不能误将定语“英雄”作为句子的主语；“得”是补语结构助词。）

至于我们开头所引的《土地》中的那个长句，句子最前面的四个“当······时候”是介词结构，充当“涌现”这个动词谓语的时间状语；副词“曾否”和介词结构“为土地”处在“涌现”之前，分别作时间状语和目的状语用（注意：“为土地”这个介词结构中的介词“为”不能看作谓语，因为介词不能单独充当谓语，即使介词后面带上名词组成介词结构也不能充当谓语）；“许许多多”后面有定语结构助词“的”，又处在宾语“遐想”之前，所以是定语。至于“遐想”之后的破折号，其作用是引起三个动宾词组作“遐想”这个宾语的复指成分（严格地说，作复指成分的应是三个“想起”之后的宾语）。

说到这里，我们就可以回答开头提出的怎样分析这个长句的问题了。简言之，秦牧同志所写的这个长句，是个比较复杂的单句，而且是个主谓句。它的整个结构就是：

[当······的时候]，\uuline{你}$\Vert$ [曾否] [为土地]\uline{涌现}过（许许多多）的\uwave{遐想}--\uwave{想起（它的）过去，（它的）未来······}?
\newpage